\documentclass[conference]{IEEEtran}
\usepackage[utf8]{inputenc}
\usepackage[T1]{fontenc}
\usepackage[spanish,es-tabla]{babel}
\spanishplainpercent
\usepackage{amsmath,amssymb,booktabs,url}
\usepackage[hidelinks]{hyperref}
\hypersetup{pdfauthor={Kevin Aguilar, Kenneth Campos y Kevin Cortés}}
\renewcommand{\IEEEkeywordsname}{Palabras clave}
\pagestyle{plain}
\raggedbottom
\newcommand{\pendiente}[1]{\par\medskip\noindent\textit{Pendiente de completar: #1.}\par\medskip}
\newcommand{\estado}{\noindent\textbf{Documento de trabajo. Validación experimental pendiente.}\par\medskip}
\newcommand{\autores}{\author{\IEEEauthorblockN{Kevin Aguilar, Kenneth Campos y Kevin Cortés}
\IEEEauthorblockA{Escuela de Ingeniería Electrónica\\Instituto Tecnológico de Costa Rica\\
EL4201 Procesamiento Electrónico de Potencia}}}
\newcommand{\referencias}{\begin{thebibliography}{9}\raggedright
\bibitem{enunciado} EL3313 y EL4201, \textit{Proyecto 4: instructivo y rúbricas}, secc. 4.3, 8.2 y anexo B, II semestre, 2026.
\bibitem{irf} Infineon, \textit{IRF540NPbF}, hoja de datos. [En línea]. \url{https://www.infineon.com/assets/row/public/documents/24/49/infineon-irf540n-datasheet-en.pdf}.
\bibitem{ir} Infineon, \textit{IR2110/IR2113: High and Low Side Driver}, hoja de datos. [En línea]. \url{https://www.infineon.com/assets/row/public/documents/24/49/infineon-ir2110-datasheet-en.pdf}.
\bibitem{acs} Allegro MicroSystems, \textit{ACS712}, hoja de datos, rev. 22. [En línea]. \url{https://www.allegromicro.com/-/media/Files/Datasheets/ACS712-Datasheet.ashx}.
\bibitem{pr} Imperix, \textit{Proportional resonant controller}, TN110, 2024. [En línea]. \url{https://imperix.com/doc/implementation/proportional-resonant-controller}.
\bibitem{plecs} Plexim, \textit{PLECS 5.0 User Manual}. [En línea]. \url{https://www.plexim.com/sites/default/files/plecsmanual-5-0.pdf}.
\bibitem{xadc} AMD, \textit{7 Series FPGAs and Zynq-7000 SoC XADC User Guide}, UG480, v1.11, 2022, cap. 2. [En línea]. \url{https://www.amd.com/content/dam/xilinx/support/documents/user_guides/ug480_7Series_XADC.pdf}.
\end{thebibliography}}

\hypersetup{pdftitle={Diseño del inversor monofásico: documento de trabajo}}
\begin{document}
\title{Diseño de un inversor monofásico con control proporcional resonante de corriente}
\autores\maketitle\estado
\begin{abstract}
Se propone un puente H para regular corriente sinusoidal de 50 Hz en una carga RL con alimentación de 12 a 24 V y referencias de 0,2, 0,5 y 1 A pico. Se documentan la topología, los componentes disponibles, el dimensionamiento, el modelo promediado, un candidato PR discreto y el plan de validación. Las simulaciones respaldan la evaluación inicial bajo hipótesis identificadas. El accionamiento, la medición y los parámetros finales deberán contrastarse experimentalmente.
\end{abstract}
\begin{IEEEkeywords}inversor monofásico, PWM unipolar, control PR, corriente.\end{IEEEkeywords}

\section{Objetivos y alcance}
El objetivo es diseñar y comprobar un lazo de corriente sinusoidal. El instructivo exige error de amplitud fundamental no mayor del 10\,\% y rizado de conmutación no mayor de 0,20 A pico a pico en los puntos nominales \cite{enunciado}. El desfase se mide y reporta, sin límite adicional. La ejecución sobre MicroBlaze V se coordina con Taller; la evaluación propia de Potencia requiere igualmente una demostración física con PR discreto.

\section{Topología, carga y selección}
El puente H utiliza cuatro MOSFET de canal N. En PWM unipolar, las ramas comparan referencias de signos opuestos con una portadora común y producen niveles diferenciales $+V_{dc}$, 0 y $-V_{dc}$. Se distinguen la frecuencia por transistor y las componentes predominantes en la tensión diferencial.

\begin{table}[ht]\caption{Componentes principales disponibles}\centering\small
\begin{tabular}{@{}cll@{}}\toprule
Cant. & Elemento & Función\\\midrule
4 & IRF540N & Interruptores\\
2 & IR2110 DIP-14 & Accionamiento por rama\\
1 & ACS712 de 5 A & Medición de corriente\\
3 & $3,3\ \Omega$, 5 W & Carga resistiva\\
1 & Secundario de transformador & Carga inductiva\\\bottomrule
\end{tabular}\end{table}

Dos resistencias forman $R_{\mathrm{ext}}=6,6\ \Omega$; la tercera permite llegar a $9,9\ \Omega$, un aumento del 50\,\%. La inserción debe conservar un camino para la corriente inductiva.

Se midieron 9,291 mH a 120 Hz y 6,984 mH a 1 kHz en el secundario, con el primario abierto. La resistencia observada del devanado es aproximadamente $0,6\ \Omega$. El modelo usa 10 mH constantes como hipótesis; no es una medición a 50 Hz ni una extrapolación validada. El primario permanecerá abierto y con terminales aislados. La corriente nominal del transformador no demuestra linealidad como inductor.

\section{Modelo y dimensionamiento}
Para $R=6,6+0,6=7,2\ \Omega$:
\begin{equation}
L\frac{di_L}{dt}+Ri_L=v_{AB},\qquad G_{vi}(s)=\frac{1}{Ls+R}.
\end{equation}
En promedio ideal, $\overline v_{AB}=mV_{dc}$. Se adopta inicialmente $|m|\leq0,85$. La amplitud necesaria para una senoide es
\begin{equation}
V_{1,pk}=I_{pk}\sqrt{R^2+(2\pi f_0L)^2}.
\end{equation}
Con L de 10 mH y 50 Hz, $|Z|\approx7,856\ \Omega$. A 1 A pico se requieren aproximadamente 7,856 V pico: $m_{pk}\approx0,655$ a 12 V y 0,327 a 24 V, antes de caídas y tiempo muerto.

Para una senoide, $I_{\mathrm{rms}}=I_{pk}/\sqrt2$. A 1 A pico, cada resistencia disipa 1,65 W y el devanado pierde aproximadamente 0,30 W por resistencia de cobre. Las resistencias de 5 W ofrecen margen nominal, sujeto a temperatura y montaje. No se incluyen pérdidas magnéticas.

\subsection{Semiconductores y pérdidas}
El IRF540N tiene tensión nominal de 100 V y Ron máxima de 44 m$\Omega$ especificada a $V_{GS}=10$ V y 25 $^\circ$C \cite{irf}. El umbral de compuerta no equivale a una tensión suficiente de accionamiento; se requieren drivers.

Con dos canales conduciendo, $P_{\mathrm{cond}}\approx2R_{\mathrm{DS(on)}}I_{\mathrm{rms}}^2$, equivalente a 0,044 W para 1 A pico y 44 m$\Omega$, sin pérdidas de diodos ni variación térmica. La estimación de conmutación se completará con
\begin{equation}
P_{\mathrm{sw}}\approx\sum_j \tfrac12 V_{\mathrm{DS},j}|I_j|
(t_{r,j}+t_{f,j})f_{\mathrm{sw}}.
\end{equation}
\pendiente{tiempos efectivos con las resistencias de compuerta elegidas, pérdidas de diodos, temperatura y margen térmico final}

\subsection{Accionamiento y frecuencia}
Se propone un IR2110 por rama, con VDD de 3,3 V y VCC fijo de 12 V, separado del bus variable \cite{ir}. El capacitor bootstrap se conecta entre VB y VS; HO se mide respecto a VS y LO respecto a COM. El esquema de una rama y su tabla de terminales se conservan en la documentación de componentes.

Se evalúan resistencias de compuerta de $22\ \Omega$, resistencias compuerta--fuente de 10 k$\Omega$, desacoplos locales y tiempo muerto de 1 $\mu$s. Se verificarán apagado y ausencia de conducción simultánea antes de energizar el bus.

Se toma 10 kHz como frecuencia inicial, con control cada 100 $\mu$s y 200 actualizaciones por ciclo de 50 Hz. La selección final debe contrastar rizado, pérdidas, resolución y ejecución. Las frecuencias digitales adicionales de 20 y 25 kHz aún requieren validación de potencia.

\section{Medición}
El ACS712 de 5 A se conecta en serie con la carga y se alimenta a 5 V regulados. Nominalmente, $v_s=2,5+0,185i_L$ \cite{acs}. El divisor propuesto de 4,7 k$\Omega$/1 k$\Omega$ produce
\begin{equation}
v_{\mathrm{ADC}}=v_s/5,7\approx0,4386+0,03246i_L.
\end{equation}
A $\pm1$ A, resulta aproximadamente 0,406--0,471 V. El XADC se usa en el rango unipolar de 0 a 1 V y 12 bits \cite{xadc}: la resolución ideal equivale a unos 7,52 mA por código. A 0,2 A pico la excursión es de unos 27 códigos respecto al cero.

La copia funcional evalúa 10 nF en el nodo ADC, con $\tau\approx8,25\ \mu$s. No es un filtro final validado. Se medirán cero, sensibilidad, signo, ruido, precisión y retardo; si el rango útil es insuficiente, se revisará la ganancia analógica.

\section{Control PR y discretización}
Se estudia la estructura de ancho de banda finito \cite{pr}:
\begin{equation}
C(s)=K_p+\frac{2K_r\omega_cs}{s^2+2\omega_cs+\omega_0^2}.
\end{equation}
El candidato usa $K_p=12$ V/A, $K_r=200$ V/A, $\omega_0=2\pi50$ rad/s y $\omega_c=2\pi2$ rad/s. La estructura concentra ganancia alrededor de la referencia. Estos parámetros son un candidato evaluado, no una sintonización final demostrada.

Se aplica Tustin preajustado a 50 Hz:
\begin{equation}
s=\kappa\frac{1-z^{-1}}{1+z^{-1}},\qquad
\kappa=\frac{\omega_0}{\tan(\omega_0T_s/2)}.
\end{equation}
Con $T_s=100\ \mu$s, la rama resonante cumple
\begin{align}
r[k]={}&b_0(e[k]-e[k-2])\nonumber\\
&-a_1r[k-1]-a_2r[k-2],
\end{align}
con $b_0=0,2509707453602995$, $a_1=-1,996504651666988$ y $a_2=0,9974902925463971$. Se calcula $v_{\mathrm{PR}}=12e+r$ y se aplica $m[k]=\mathrm{sat}_{0,85}(v_{\mathrm{PR}}[k-1]/12)$, con normalización fija y una muestra de retardo.

Para la planta con retención, $A=e^{-RT_s/L}$ y $B=(1-A)/R$. Con $q=z^{-1}$ y $C(q)=N_C(q)/D(q)$, el polinomio cerrado es
\begin{equation}
(1-Aq)D(q)+(V_{dc}/12)Bq^2N_C(q).
\end{equation}
La recursión de Schur confirma estabilidad en doce combinaciones de R total de 7,2/10,5 $\Omega$, L de 6,984/9,291/10 mH y bus de 12/24 V. Esta comprobación lineal excluye medición, tiempo muerto y saturación.

\pendiente{procedimiento y justificación de sintonización final, márgenes, retardos medidos, saturación y anti-windup}

\section{Plan de validación}
Primero se comprobarán auxiliares, apagado, bootstrap, compuertas y tiempo muerto; se documentarán límite de corriente, fusible y desconexión. Se calibrará el ACS712 con una medida independiente y se comprobará la carga.

Se realizarán por separado: (i) las tres referencias nominales; (ii) cambios de 0,2 a 0,5 A y de 0,5 a 1 A pico; (iii) cambio resistivo de al menos 25\,\% a 0,5 A pico. Cada registro identificará referencia, corriente, tensión de carga, PWM y parámetros reales. La fundamental se extraerá en al menos cinco períodos completos. El rizado local se separará de la componente de baja frecuencia.

\section{Estado del diseño}
Se dispone de componentes principales, mediciones preliminares y modelos reproducibles. Falta completar el esquema del montaje, caracterizar la medición, cerrar pérdidas y controlador y verificar físicamente los criterios. Los datos experimentales permitirán actualizar el documento de diseño y sustentar el informe final.
\referencias
\end{document}
